\section*{Chapitre \ref{ch:b1:observing-the-sky} --- Télescopes, détecteurs et photométrie}

\begin{solution}{exo:b1:observing-the-sky:1}
Une différence de 5 magnitudes est un facteur 100~: l'étoile de magnitude 1.0
envoie 100 fois plus de \omterm{def:b1:observing-the-sky:flux}{flux}. Rigel et Polaris~: $10^{0.4\times(2.02 - 0.13)} =
10^{0.756} = 5.70$~; Rigel, de plus petite magnitude, est la plus brillante.
\end{solution}

\begin{solution}{exo:b1:observing-the-sky:2}
$d = 1/0.19495 = \qty{5.13}{\parsec}$, soit $5.1295\times206265 =
\qty{1.058e6}{\au}$ et $5.1295\times3.2616 = \qty{16.7}{\lightyear}$.
\end{solution}

\begin{solution}{exo:b1:observing-the-sky:3}
$1.22\times\num{550e-9}/2.4 = \qty{2.80e-7}{rad} = \ang{;;0.058}$, soit environ
sept fois mieux qu'un \omterm{def:b1:observing-the-sky:seeing}{seeing} de \ang{;;0.40}~: dans l'espace, le télescope
atteint sa limite de diffraction. Pour \qty{6.5}{m} à \qty{2.0}{\micro\metre}~:
$1.22\times\num{2.0e-6}/6.5 = \qty{3.75e-7}{rad} = \ang{;;0.077}$.
\end{solution}

\begin{solution}{exo:b1:observing-the-sky:4}
$B - V$ croît de Rigel ($-0.03$) à Altaïr (0.22), Polaris (0.60) et
Bételgeuse (1.85)~: la plus chaude est Rigel, puis Altaïr, Polaris et
Bételgeuse. La distance affaiblit les deux bandes du même facteur, donc $B - V$
n'en est pas modifié (tant qu'aucune poussière ne rougit la lumière en chemin).
\end{solution}

\begin{solution}{exo:b1:observing-the-sky:5}
$X_1 = \sec\ang{30} = 1.1547$, $X_2 = \sec\ang{60} = 2.0000$.
$k = (6.590 - 6.412)/(2.0000 - 1.1547) = 0.178/0.8453 = \num{0.211}$ par
\omterm{def:b1:observing-the-sky:extinction}{masse d'air}, et $m_0 = 6.412 - 0.2106\times1.1547 = 6.169$.
\end{solution}

\begin{solution}{exo:b1:observing-the-sky:6}
$\mathrm{SNR} = \sqrt{400} = 20$. Le \omterm{def:b1:observing-the-sky:detector}{rapport signal sur bruit} croît comme
$\sqrt t$~: pour 100,
$t = \qty{10}{s}\times(100/20)^2 = \qty{250}{s}$. Une erreur de \num{0.01}
magnitude exige $\mathrm{SNR} = 1.086/0.01 = 108.6$, donc
$t = 10\times(108.6/20)^2 = \qty{295}{s}$.
\end{solution}

\begin{solution}{exo:b1:observing-the-sky:7}
La distance est $d = 1/0.00378 = \qty{264.6}{\parsec}$, et la \omterm{def:b1:observing-the-sky:absolute}{magnitude absolue}
$M = 0.13 - 5\log_{10}26.46 = 0.13 - 7.11 = -6.98$. $L/L_\odot = 10^{0.4\times(4.83 + 6.98)} =
10^{4.724} = \num{5.3e4}$ dans V.
\end{solution}

\begin{solution}{exo:b1:observing-the-sky:8}
$\pi\times1.2^2 = \qty{4.52}{m^2}$, et $25/4.52 = 5.5$~: le télescope Webb
recueille le même nombre de photons en une pose 5.5 fois plus courte. Un
disque plein de \qty{6.5}{m} de diamètre aurait $\pi\times3.25^2 = \qty{33.2}{m^2}$~;
le miroir réel en a $25/33.2 = \qty{75}{\percent}$, car 18 hexagones laissent
découvert le bord de ce disque et les interstices entre segments et l'ouverture
centrale ne recueillent rien.
\end{solution}

\begin{solution}{exo:b1:observing-the-sky:9}
(a) $\qty{10}{\parsec} = \qty{3.0857e17}{m}$, et $F = \num{3.828e26}/(4\pi
\times(\num{3.0857e17})^2) = \num{3.828e26}/\num{1.1965e36} =
\qty{3.20e-10}{W.m^{-2}}$, $1361/\num{3.20e-10} = \num{4.3e12}$ fois moins qu'à
la Terre.
(b) $5\log_{10}(d/\qty{10}{\parsec}) = -26.74 - 4.83 = -31.57$, donc
$d = \qty{10}{\parsec}\times10^{-6.314} = \qty{4.853e-6}{\parsec}$, soit
$\num{4.853e-6}\times206265 = \qty{1.001}{\au}$.
\end{solution}

\begin{solution}{exo:b1:observing-the-sky:10}
Au-dessus de toute l'atmosphère le \omterm{def:b1:observing-the-sky:flux}{flux} est $F_0$~; au-dessous, $F = F_0e^{-\tau_0X}$.
Alors $m - m_0 = -2.5\log_{10}(F/F_0) = 2.5\log_{10}e\;\tau_0X$, ce qui est $kX$
avec $k = 2.5\log_{10}e\;\tau_0 = 1.0857\,\tau_0$. Pour $k = 0.160$~:
$\tau_0 = 0.160/1.0857 = 0.147$. Au \omterm{def:b1:spherical-astronomy:celestial-sphere}{zénith}, $e^{-0.147} = 0.863$ de la lumière
arrive, soit \qty{86}{\percent}~; pour $X = 2.31$, $e^{-0.147\times2.31} =
e^{-0.340} = 0.712$, soit \qty{71}{\percent}. Pour $z = \ang{85}$, $\sec z =
11.5$~; mais le modèle plan-parallèle rend le trajet infini à l'horizon, alors
que le trajet à travers une atmosphère courbe reste fini, et la réfraction
courbe aussi le rayon~: près de l'horizon, la formule surestime la \omterm{def:b1:observing-the-sky:extinction}{masse d'air}
et ses erreurs croissent vite.
\end{solution}

\begin{solution}{exo:b1:observing-the-sky:11}
(a) $\mathrm{SNR} = Rt/\sqrt{(R + B)t}$, donc $t = \mathrm{SNR}^2(R + B)/R^2
= 100\times2050/2500 = \qty{82}{s}$. (b) Diviser le \omterm{def:b1:observing-the-sky:seeing}{seeing} par deux divise par 4
l'aire de l'image, donc $B$~: $B = 500$ et $t = 100\times550/2500 =
\qty{22}{s}$. (c) Pour $B \gg R$, $t \approx \mathrm{SNR}^2B/R^2$. Le taux de la
source $R$ croît comme $D^2$~; le taux du ciel $B$ croît comme la surface
collectrice $D^2$ fois l'aire de ciel de l'image, $\theta^2$. Donc $t \propto
\theta^2D^2/D^4 = \theta^2/D^2$~: diviser le \omterm{def:b1:observing-the-sky:seeing}{seeing} par deux vaut autant que
doubler le diamètre.
\end{solution}

\begin{solution}{exo:b1:observing-the-sky:12}
(a) $N = 1200\times100 = \num{120000}$ traits~; au premier ordre, le pouvoir de
résolution est \num{1.2e5}, et $\Delta\lambda = 550/(\num{1.2e5}) =
\qty{4.6e-3}{nm}$. (b) $v = c\,\Delta\lambda/\lambda = c/(\num{1.2e5}) =
\qty{2.5}{km.s^{-1}}$. (c) $\Delta\lambda = 550\times29.78/\num{299792} =
\qty{0.055}{nm}$, soit environ 12 éléments de résolution. L'observateur doit
calculer, pour la date et l'heure de l'observation, la composante de la vitesse
orbitale de la Terre vers l'étoile, et la soustraire, pour obtenir la vitesse de
l'étoile par rapport au Soleil.
\end{solution}

\begin{solution}{pb:b1:observing-the-sky:1}
\textbf{1.} $d = 1/0.37921 = \qty{2.637}{\parsec}$.

\textbf{2.} $2.637\times206265 = \qty{5.44e5}{\au}$;
$2.637\times\num{3.0857e16} = \qty{8.14e16}{m}$.

\textbf{3.} $2.637\times3.2616 = \qty{8.60}{\lightyear}$~: la lumière a quitté
Sirius il y a 8.6 ans.

\textbf{4.} La direction oscille de $2p = \ang{;;0.758}$, soit environ deux fois
le \omterm{def:b1:observing-the-sky:seeing}{seeing}. L'oscillation se mesure à partir de la position du centre de l'image,
non de son bord~: en moyenne sur de nombreux photons, le centre est trouvé à une
petite fraction de la largeur de l'image, d'autant plus précisément que l'on
compte plus de photons.

\textbf{5.} $\Delta d/d = \Delta p/p = 0.001/0.37921 = \qty{0.26}{\percent}$,
soit $8.60\times0.0026 = \qty{0.023}{\lightyear}$.

\textbf{6.} $10^{0.4\times(0.03 + 1.46)} = 10^{0.596} = 3.94$.

\textbf{7.} $10^{0.4\times(26.74 - 1.46)} = 10^{10.112} = \num{1.29e10}$.

\textbf{8.} $10^{0.4\times(8.44 + 1.46)} = 10^{3.96} = \num{9120}$.

\textbf{9.} $z = \ang{90} - \ang{82.09} = \ang{7.91}$, $X = \sec\ang{7.91} =
1.0096$~; extinction $0.15\times1.0096 = 0.151$ magnitude~; magnitude mesurée
$m = -1.46 + 0.151 = -1.31$.

\textbf{10.} $z = \ang{60}$, $X = 2.000$~; extinction 0.30 magnitude,
$m = -1.16$. La fraction transmise est $10^{-0.4\times0.30} = 0.759$~:
\qty{24}{\percent} de la lumière est perdue.

\textbf{11.} Au lever, $z = \ang{90}$ et $\sec z$ est infini, alors que le
trajet réel à travers l'atmosphère courbe est long mais fini~; la réfraction
courbe aussi le rayon. Près de l'horizon, l'extinction est grande et change vite,
et la correction plan-parallèle est fausse.

\textbf{12.} $\num{3.0e4}\times10^{0.4\times(10.0 - 8.44)} = \num{3.0e4}
\times10^{0.624} = \num{3.0e4}\times4.21 = \num{1.26e5}$ photoélectrons par
seconde.

\textbf{13.} Un \omterm{def:b1:observing-the-sky:detector}{rapport signal sur bruit} de 1000 exige $N = 1000^2 = \num{1.0e6}$
photoélectrons~: $t = \num{1.0e6}/\num{1.26e5} = \qty{7.9}{s}$. L'erreur sur la
magnitude est $1.086/1000 = \num{0.0011}$.

\textbf{14.} $\num{3.0e4}\times10^{0.4\times(10.0 + 1.46)} = \num{3.0e4}
\times10^{4.584} = \num{1.15e9}$ photoélectrons par seconde.

\textbf{15.} $0.20\times\num{1.15e9} = \num{2.30e8}$ par seconde dans le pixel
le plus brillant, qui se remplit en $\num{1.0e5}/\num{2.30e8} =
\qty{4.3e-4}{s}$, soit moins d'une demi-milliseconde.

\textbf{16.} Les ailes de Sirius donnent $\num{1.15e9}\times10^{-4} =
\num{1.15e5}$ par seconde~; Sirius~B donne $0.20\times\num{1.26e5} =
\num{2.52e4}$ par seconde. L'éblouissement de Sirius est 4.6 fois plus brillant
que le pic de Sirius~B.

\textbf{17.} Le pixel reçoit $\num{1.15e5} + \num{2.52e4} = \num{1.40e5}$
par seconde et se remplit en $\num{1.0e5}/\num{1.40e5} = \qty{0.71}{s}$. En
\qty{0.5}{s}~: signal $\num{2.52e4}\times0.5 = \num{1.26e4}$, total
$\num{7.01e4}$, bruit $\sqrt{\num{7.01e4}} = 265$, donc $\mathrm{SNR} =
\num{1.26e4}/265 = 48$. Sirius~B est détectable en ce qui concerne le bruit de
photons~; la difficulté tient à ce qu'il faut connaître et soustraire
l'éblouissement avec une précision bien meilleure qu'un cinquième de sa valeur,
alors que le cœur de Sirius sature en moins d'une milliseconde.

\textbf{18.} $m - M = 5\log_{10}(2.637/10) = 5\times(-0.579) = -2.89$.

\textbf{19.} $M = -1.46 + 2.89 = 1.43$.

\textbf{20.} $d = 1/0.13023 = \qty{7.679}{\parsec}$~;
$m - M = 5\log_{10}0.7679 = -0.57$, donc $M = 0.03 + 0.57 = 0.60$.

\textbf{21.} Véga, de plus petite \omterm{def:b1:observing-the-sky:absolute}{magnitude absolue}, d'un facteur
$10^{0.4\times(1.43 - 0.60)} = 10^{0.332} = 2.15$. Sirius paraît plus brillant
seulement parce qu'il est plus proche.

\textbf{22.} $m = M_\odot + (m - M) = 4.83 - 2.89 = 1.94$~: le Soleil serait une
étoile assez brillante à l'œil nu.

\textbf{23.} $5\log_{10}(d/\qty{10}{\parsec}) = 6.0 - 1.43 = 4.57$, donc
$d = 10\times10^{0.914} = \qty{82.0}{\parsec}$, soit environ \qty{267}{\lightyear}.

\textbf{24.} $L/L_\odot = 10^{0.4\times(4.83 - 1.43)} = 10^{1.36} = 22.9$~:
dans la bande V, Sirius éclipse \textbf{environ 23 soleils}.
\end{solution}
